\documentclass{article}
\usepackage{../../shared/omnicourse}
\usepackage{graphicx}
\input{../../shared/meta.tex}

\renewcommand{\docshortname}{UI Prototypes}
\hypersetup{pdftitle={OmniCourse -- UI Prototypes v2.0}}

% Screenshots are captured separately from docs/analysis-and-design/ui-prototypes (English, light theme).
\providecommand{\shotdir}{shots}
\graphicspath{{\shotdir/}}
% \shot[max height]{file}{caption}
\newcommand{\shot}[3][0.62\textheight]{%
  \begin{figure}[H]\centering
    {\setlength{\fboxsep}{0pt}\color{ocRule}\fbox{\includegraphics[width=\textwidth,height=#1,keepaspectratio]{#2}}}
    \caption{#3}
  \end{figure}}
% step-to-screen table
\newcommand{\stepmap}[1]{\par{\small\begin{xltabular}{\dimexpr\textwidth-2\tabcolsep\relax}{@{}p{2.4cm}Y@{}}
  \hcell{Step} & \hcell{Where it is on the screen}\\\hdrline\endhead #1\end{xltabular}}\par}

\begin{document}

\makecover{UI Prototypes}{2.0}{1 November 2026}{ESE-PA2-UI-2.0}{Project Assignment 2 (PA2)}

\pagenumbering{roman}
\thispagestyle{empty}

\section*{Revision history}
\vspace{-0.3em}
{\small
\begin{xltabular}{\dimexpr\textwidth-2\tabcolsep\relax}{@{}p{2.3cm}p{1.4cm}Yp{3.6cm}@{}}
\hcell{Date} & \hcell{Version} & \hcell{Description} & \hcell{Author}\\\hdrline
\endhead
18/10/2026 & 1.0 & UI sketches of the main screens for PA1. & Lê Hữu Hoà\\
31/10/2026 & 1.5 & First clickable HTML prototypes of every screen (PR \#44). & Lê Hữu Hoà\\
1/11/2026 & 2.0 & Redesigned prototype set: 16 pages and 27 screens and states, English and Vietnamese, light and dark theme,
collapsible sidebar, planner with one row per term and live rule checks, grades with per-part bars (PR \#48). Every screen
mapped to the steps of its use case. & Lê Hữu Hoà\\
\end{xltabular}}

\vspace{1.2em}
{\hypersetup{linkcolor=ocNavy}\tableofcontents}
\clearpage
\pagenumbering{arabic}

%% =====================================================================
\section{Introduction}
The PA2 brief asks for the UI of at least three key use cases, matching the use-case specifications. This document
covers twelve of the nineteen use cases. For each one it shows the screens and a table that ties each step of the basic
flow, and the main alternative flows, to what the student sees. Step numbers refer to the Use-Case Model and
Specification v1.0 (ESE-PA2-UCS-1.0).

\subsection{How to open the prototype}
The prototype is a set of static HTML pages, included in this package in the folder \texttt{ui-prototypes}. Open
\texttt{index.html} in a browser; nothing needs installing. The pages share synthetic data: the student
\emph{Nguyễn Minh Phương}, class APCS 24A01, student ID 24125369, all made up. Press \textbf{G} on any page to jump to
any screen or state. The globe button switches between Vietnamese and English; the screenshots here are in English.

\subsection{Design direction}
\begin{itemize}
  \item \textbf{One colour and one glyph per requirement group.} Core computing is blue \texttt{</>}, mathematics orange
        $\Sigma$, physics teal, English yellow, political theory rose, general education green, electives violet and the
        thesis graphite. The glyph keeps the meaning when colour is not seen.
  \item \textbf{Course tiles.} A solid tile is passed, a tile with an inner ring is in progress, a white tile with a side
        bar is planned, and a dashed tile is an empty slot.
  \item \textbf{Calm screens.} A light grey page, white cards, colour kept for groups and for status. Figures read heavier
        than the text around them, and units are written in full (\emph{22 credits}, never \emph{22 cr}).
  \item \textbf{The rule engine is visible.} Wherever a rule decides something, the screen says which rule and cites
        the regulation clause, so the student can check it.
\end{itemize}

\subsection{Screens and use cases}
{\small
\begin{xltabular}{\dimexpr\textwidth-2\tabcolsep\relax}{@{}p{3.6cm}Yp{2.6cm}@{}}
\hcell{Page} & \hcell{What it shows} & \hcell{Use cases}\\\hdrline
\endhead
\texttt{index.html} & Landing page for visitors & UC-01, UC-13\\
\texttt{sign-up.html}, \texttt{sign-in.html} & Account creation with consent; sign-in & UC-01, UC-13\\
\texttt{onboarding.html} & Program, cohort, target term, credits per term, transcript & UC-01\\
\texttt{import.html} & Upload and review of the transcript & UC-02\\
\texttt{dashboard.html} & Degree audit & UC-03\\
\texttt{audit.html} & Grades and what each course still needs & UC-15\\
\texttt{explorer.html} & Curriculum list and prerequisite graph & UC-04\\
\texttt{planner.html} & Plan, goal box, what-if, share & UC-05, UC-06, UC-07, UC-09\\
\texttt{shared.html} & Read-only shared plan; dead link page & UC-16\\
\texttt{advisor.html} & Advisor chat & UC-08\\
\texttt{guides.html} & School-task guides & UC-10\\
\texttt{inbox.html} & Notifications & UC-12\\
\texttt{settings.html} & Account, privacy, sessions, consent & UC-14\\
\texttt{admin.html} & Curriculum administration & UC-11, UC-17\\
\end{xltabular}}

%% =====================================================================
\section{UC-01 Sign up and onboard}
\shot[0.38\textheight]{sign-up}{Sign-up: email, password with strength meter, display name, consent with a link to the
privacy notice.}
\shot[0.38\textheight]{onboarding}{Onboarding, first of five questions: the program, from the programs in the curriculum
data.}
\stepmap{
Steps 1--3 & \textbf{Create account} form on \texttt{sign-up.html}; \textbf{Continue with Google} above it (A4).\\
Step 4, A2 & Each field shows its error under it, in words; the password strength meter updates while typing.\\
Steps 6--9 & \texttt{onboarding.html}: program, cohort, target graduation term, credits per term; \textbf{Next} on each,
\textbf{Skip for now} keeps onboarding resumable (A7).\\
Step 10, A5 & The last question offers to import the transcript or do it later.\\
Step 11 & \textbf{Finish} opens the dashboard.\\
}

%% =====================================================================
\section{UC-02 Import transcript}
\shot[0.3\textheight]{import}{Upload: the drop zone, the note that the PDF is read on the server and deleted, and the link
to manual entry.}
\shot[0.55\textheight]{import-review}{Review: every row with its code, credits, grade and term. One row could not be read
and blocks the save until it is fixed.}
\stepmap{
Steps 1--3 & Drop zone on \texttt{import.html}, with \textbf{How to get the PDF from the portal} and the privacy note.\\
Steps 4--6 & Review table: rows in progress are marked \textbf{In progress} (A3); a row not in the curriculum is marked
in words; each row can be edited or removed, \textbf{Add row} adds one (A4).\\
A2 & The unreadable row is highlighted, and the header says \emph{1 row to fix}; \textbf{Save transcript} stays disabled.\\
Steps 7--9 & \textbf{Save transcript}, then the dashboard opens.\\
}

%% =====================================================================
\section{UC-03 View degree audit}
\shot{dashboard}{Dashboard: credits by requirement group, GPA with its change, expected graduation, credits left, one bar
per group, this term's courses, \emph{Needs attention} and the by-term strip.}
\stepmap{
Step 3 & Top row: cumulative GPA with the change from last term, credits passed of 167, credits left, expected graduation
term.\\
Step 4 & Donut chart of credits: each group's passed credits, then in progress (hatched), then still to take in light grey.
Hovering a group lights it in the donut and in the list.\\
Steps 5, 7, 8 & \emph{Requirement groups}: one bar per group; clicking a bar lists passed, in-progress and still-needed
courses.\\
Step 6 & \emph{Needs attention}: plan issues, a missing language certificate, an empty elective group, the registration
window, each with a link to act.\\
A2 & Hovering a GPA point or a past term previews the page as it stood at the end of that term; the arrows and the term
picker in the header keep it there.\\
}

%% =====================================================================
\section{UC-15 View grades and what each course still needs}
\shot[0.7\textheight]{audit}{Grades: three figures, the term averages with the pass (5) and excellent (8) lines, then each
course with its graded parts stacked on a 0--10 bar.}
\stepmap{
Step 2 & Cumulative GPA, credits passed, best term; the line chart with lines at 5 and 8; a failed course makes a red
point.\\
Step 3 & Current term, planned terms, then past terms; each course has a stacked bar, thin cutoffs at 5 and 8, the grade
and a status pill (\textbf{In progress}, \textbf{Planned}, \textbf{Passed}, \textbf{Failed}).\\
Step 4, A4, A5 & Under an unfinished course: \emph{Final exam needs 1.4 to pass, 7.4 for excellent}, or \emph{Average 5.0
to pass, 8.0 for excellent} when no part is graded.\\
Steps 5--6 & Clicking a course opens each graded part with its weight and, for a retake, every sitting.\\
A1, A2 & Group chips and the search box above the list; \emph{No course matches} when the filter is empty.\\
}

%% =====================================================================
\section{UC-04 Explore curriculum}
\shot[0.5\textheight]{explorer}{Curriculum list: in each group, courses in progress, then available, then locked with the
missing prerequisites, then passed.}
\shot[0.4\textheight]{graph}{Prerequisite graph (A1): clicking a node lights its whole chain.}
\stepmap{
Steps 1--3 & Group sections with the order in progress, available, locked (padlock and \emph{needs \ldots}), passed.\\
Steps 4--5 & Search box; typing without diacritics matches Vietnamese names.\\
Steps 6--7 & Clicking a course opens the drawer with credits, prerequisites, what it unlocks and the terms offered.\\
A1 & \textbf{Prerequisite graph} tab; drag or scroll to pan.\\
}

%% =====================================================================
\section{UC-05 Build and edit a plan}
\shot[0.7\textheight]{planner}{Planner: the status strip, one row per term with a credit meter, the course pool, and the
Issues panel.}
\stepmap{
Step 2 & Status strip: estimated graduation (or \emph{Not valid} with the count) and the coverage of each group. Rows:
\emph{Completed} (one row of codes), the current term, then each future term with its load and a meter marking the
program's minimum and cap. \emph{Issues} panel on the right.\\
Steps 3--4 & While a course is dragged, terms that cannot take it fade and refuse the drop, the target row previews its new
load, every row turns green or red with the reason, and the panel previews new and resolved issues.\\
Steps 5--6 & The drop updates meters, coverage, the estimated graduation term and the panel.\\
A1 & Space lifts a tile, up and down arrows move it and skip refused terms, Enter drops.\\
A2 & Each issue has its clause and, when one exists, \textbf{Apply} for the fix; \textbf{Apply every suggestion} at the
bottom.\\
A3, A6 & \textbf{Undo} in the header; the plan menu creates, duplicates and deletes plans (delete asks first).\\
}

%% =====================================================================
\section{UC-07 Generate a plan}
\shot[0.45\textheight]{planner-goal}{Goal box: a sentence fills the planning-goal form; the changed fields are marked
before the student presses Plan.}
\stepmap{
Step 2 & The form under the goal box: target term, maximum credits, internship term, courses to avoid; a note says the
text is sent to an AI provider.\\
Steps 3--5 & \textbf{Fill the form} fills the fields and marks each one that changed, in words and colour.\\
Step 6 & \textbf{Re-plan from the form}; the rule engine checks the result, which shows in the rows and the Issues panel.\\
A1, A2 & The form works without the goal box; when AI is unavailable the box says so and the form stays usable.\\
}

%% =====================================================================
\section{UC-06 Simulate a what-if scenario}
\shot[0.45\textheight]{planner-whatif}{What-if: the scenario runs on a copy; the panel shows the graduation term before and
after and the courses that move.}
\stepmap{
Steps 1--3 & \textbf{What if} in the planner header opens the panel with \textbf{Fail a course}, the course picker and
the other scenarios.\\
Steps 4--5 & \emph{Current} and \emph{After the change} graduation terms side by side, with the list of terms and courses
that move.\\
Steps 6--7, A5 & \textbf{Save as a new plan} keeps the original; \textbf{Close} discards the copy.\\
}

%% =====================================================================
\section{UC-09 Share or export a plan and UC-16 View a shared plan}
\shot[0.45\textheight]{planner-share}{Share dialog: what the advisor will and will not see, the option to show the
student's name, the warning, Create link and Download PDF.}
\shot[0.6\textheight]{shared}{The advisor's read-only view: the same rows, meters and violations as the planner, with nothing
to drag or fix.}
\stepmap{
UC-09 steps 1--2 & \textbf{Share} in the planner; the dialog lists what a reader sees and what stays private (grades,
email, student ID, transcript, other plans) and \textbf{Show my name}.\\
UC-09 steps 3--4 & \textbf{Create link}, then \textbf{Copy}; \textbf{Stop sharing} once a link exists (A2, A3).\\
UC-09 A1 & \textbf{Download PDF} in the same dialog.\\
UC-16 step 3 & \texttt{shared.html}: banner \emph{Shared plan, read-only, last changed \ldots}, status strip, term rows,
\emph{Violations} panel with \emph{The rule engine checked this when the page opened}.\\
UC-16 A1 & \texttt{shared.html\#dead}: \emph{This plan is not shared}, the same page for every dead or wrong link.\\
UC-16 A2 & \textbf{Export PDF} prints the view.\\
}

%% =====================================================================
\section{UC-08 Ask the advisor chat}
\shot[0.42\textheight]{advisor-consent}{First use: what is sent to the AI provider and what never is, with I agree and Not
now.}
\shot[0.5\textheight]{advisor}{An answer: the steps the chat took, the rule engine result for this student, the text, and
the cited clauses.}
\stepmap{
Step 2, A1 & Consent dialog: sent (program, cohort, course codes, grades, terms), never sent (name, email, student ID);
\textbf{I agree} or \textbf{Only regulations}.\\
Step 3 & Empty conversation with example questions; the note that the cited clause is the authority.\\
Steps 4--7 & The question, then the steps (\emph{Search the regulations}, \emph{Read the record}, \emph{Rule engine:
\ldots}), then the answer as it is written.\\
Step 8 & The rule engine card with the result for this student, and clause chips such as \emph{Article 7.2} that open the
clause text.\\
A2, A4 & \texttt{advisor.html\#unavailable} and \texttt{\#limit}: the chat says why it is unavailable and keeps the
question.\\
}

%% =====================================================================
\section{UC-10 Follow a school-task guide}
\shot[0.45\textheight]{guides}{Guides: search, the next deadlines, and the official portals with direct links.}
\shot[0.6\textheight]{guide}{One guide: steps quoting the portal's buttons, what to prepare, the credit limit from the rule
engine, the course codes from the student's plan, and the date the guide was checked.}
\stepmap{
Steps 1--3 & Search box (\emph{What do you need to do?}) and the task cards; the portal list below with direct links.\\
Steps 4--5 & Guide page: numbered steps, \emph{When}, \emph{Account}, notes, version and last-checked date.\\
Step 6 & \textbf{Open the student portal} opens the deep link in a new tab.\\
A3 & \emph{Your credit limit: 10--22 credits per term, computed by the rule engine}, with its clause.\\
A4 & \emph{Course codes from your plan}, ready to copy; OmniCourse never registers.\\
A6 & \textbf{Report an error} at the bottom of the guide.\\
}

%% =====================================================================
\section{UC-12 Receive notifications}
\shot[0.42\textheight]{inbox}{Notifications: plan issues, the registration window, a transcript change, each with its
action; filters on the right.}
\stepmap{
Steps 4--5 & The bell in the sidebar shows the unread count and opens a short popover.\\
Steps 6--7 & \texttt{inbox.html}: newest first, each with a button to act (open the plan, see the guide).\\
A1, A2 & \emph{Transcript changed} and \emph{Plan has 2 issues} entries.\\
A3, A4 & \textbf{Mark all read}; filters by kind on the right.\\
}

%% =====================================================================
\section{UC-14 Manage account and privacy}
\shot[0.6\textheight]{settings}{Settings: profile and planning goal, privacy and AI consent, notifications, sessions, export
and deletion.}
\stepmap{
A2 & \emph{Program and goals}: program, target term, credits per term, internship.\\
A3 & \emph{Privacy and AI}: the advisor-chat consent switch and the calls used today.\\
Steps 3--6 & \emph{Sessions}: each device with \textbf{Sign out}.\\
A4, A5 & \textbf{Export my data}; \textbf{Delete account} with the 30-day grace period and a confirmation.\\
A6 & Language and theme.\\
}

%% =====================================================================
\section{UC-11 Maintain curriculum and UC-17 Maintain the course catalog}
\shot[0.5\textheight]{admin}{Curriculum administration: program and cohort, the requirement groups with their credits, and
the courses of a group.}
\stepmap{
UC-11 steps 1--2 & Program and cohort pickers; the groups with required credits; \textbf{Add cohort} (A1).\\
UC-11 steps 3--6 & Editing a group or adding a course to it; code format and prerequisite-cycle checks show under the
field (A2, A3).\\
UC-17 & The \emph{Course catalog} tab; changing a shared course shows every curriculum it reaches before saving.\\
}

\end{document}
